% Generated draft: requires mathematical and source review.
\documentclass[11pt]{article}
\usepackage[margin=25mm]{geometry}
\usepackage{fontspec}
\setmainfont{Latin Modern Roman}
\usepackage{amsmath,amssymb,mathtools,unicode-math}
\setmathfont{Latin Modern Math}
\usepackage{enumitem,needspace,xurl,hyperref}
\setlength{\parindent}{0pt}
\setlength{\parskip}{5pt}
\setlength{\emergencystretch}{3em}
\newcommand{\N}{\mathbb{N}}
\newcommand{\Z}{\mathbb{Z}}
\newcommand{\Q}{\mathbb{Q}}
\newcommand{\R}{\mathbb{R}}
\newcommand{\C}{\mathbb{C}}
\newcommand{\F}{\mathbb{F}}
\newcommand{\A}{\mathbb{A}}
\newcommand{\PP}{\mathbb{P}}
\newcommand{\E}{\mathbb{E}}
\newcommand{\eps}{\varepsilon}
\newcommand{\dd}{\,\mathrm{d}}
\begin{document}
\section*{UnsolvedMath ID: 1900016}
\textbf{Source problem code:} catalogue entry
\subsection{Definitions and mathematical statement}
Let $\mathcal S$ denote the intended collection of faces of sails in the source's continued-fraction setting, let $\cong_{\Z}$ denote its integer-congruence relation, and let $\operatorname{freq}(F)$ be the intended frequency of a face or face class. Subject to these missing definitions, the literal assertion is
\[
\forall C\in\Z_{>0},\qquad
\#\{[F]_{\cong_{\Z}}:F\in\mathcal S,\ \operatorname{freq}(F)>C\}<\infty.
\]
Thus for each positive integer threshold, only finitely many distinct integer-congruence types have frequency strictly above that threshold. There is no claimed uniform bound independent of $C$. If frequency is attached to a class, it must be invariant under the intended congruence; if it is attached to individual occurrences, the relation between occurrence frequencies and the class predicate must be specified.

The source does not specify the dimension, sail class, face dimensions, invariant measure, normalization, or limiting procedure defining frequency. In a common lattice convention, integer congruence means equivalence under an affine lattice automorphism $x\mapsto Ux+v$ with $U\in\operatorname{GL}(n,\Z)$ and $v\in\Z^n$, but the extract does not say whether this ambient affine convention, a linear convention, or an intrinsic face-lattice convention is intended.

The threshold $C$ is retained exactly as a positive integer with the strict inequality $\operatorname{freq}(F)>C$. If frequencies were probabilities normalized to lie in $[0,1]$, no such class could occur for these thresholds. A nontrivial intended statement would therefore require a different normalization or a clarification of the threshold. In particular, the record does not authorize replacing $C$ by $1/C$, or replacing positive integer thresholds by arbitrary positive real thresholds.

The supplied moderate-confidence summary describes a related classical two-dimensional finiteness result using Gauss-map or sail statistics, and labels the higher-dimensional question open, with links to adjacent source problems. Exact definitions and citations establishing those claims are absent, so they remain attributed context rather than a verified resolution of the literal statement.
\subsection{Short English statement}
The literal claim says that above each positive integer frequency threshold there are only finitely many integer-congruence types of sail faces. The frequency measure and congruence convention are missing. The integer threshold would be vacuous for probability-normalized frequencies, so its normalization needs review before a nontrivial conjecture can be fixed.
\subsection{Sources}
\begin{itemize}
\item Local source registry: \url{data/unsolvedmath/problems.json}, permanent UnsolvedMath ID $1900016$, source code catalogue entry, including the supplied OPEN-TRIAGE summary. Dataset attribution: Ulam AI and UnsolvedMath Contributors; original-source attribution and rights are retained.
\item Karpenkov, \emph{Open problems in geometry of continued fractions} (2017), Problem 16: \url{https://arxiv.org/abs/1712.01450}. The supplied metadata reports source-TeX extraction and July 29, 2026 access, with current-status and rights review flags.
\item The summary refers generically to Karpenkov's frequency papers and Arnold's work for a classical two-dimensional analogue, but supplies no full citations or precise theorem. It links the higher-dimensional question to registry source codes AMR-018-0015 and AMR-018-0017; their full statements are not part of this record.
\end{itemize}
\end{document}
